\documentclass{situationist-preprint}

\title{Indifference and Meaning}
\msid{EXI-06-b}
\msfamily{EXI-06}
\msdiscipline{Meaning, Death and the Absurd}
\mstype{Companion manuscript}
\mskeywords{cosmic indifference; meaning in life; Russell; objective worth; despair}

\begin{document}
\maketitle

\begin{abstract}
EXI-06-a argues that the universe does not care about us, that it is indifferent, and that indifference is not hostility. It is often thought that a universe indifferent to us leaves our lives without meaning, or with a meaning we must assert in defiance of it. We argue that neither follows. Meaning, on the account of EXI-01, consists in engagement with what has worth independently of anyone's attitudes, and the universe's want of attitudes toward us takes nothing from that worth. Since indifference is not hostility, it supplies no ground for despair and no adversary for defiance. Cosmic indifference is irrelevant to meaning in life, though not to every question about meaning.
\end{abstract}

\section{The inference}

Russell's \enquote{A Free Man's Worship} pictures a world of blind forces, careless of human good, within which human ideals must be maintained by an act of proud defiance \citep{Russell1903}. The essay draws two conclusions from the indifference of the universe: that our ideals receive no support from it, and that our lives take on dignity only through resistance to it. A more pessimistic reading draws a third: that a universe which does not care about us leaves our lives without significance.

EXI-06-a establishes that the universe is indifferent: it is not responsive to our interests (Definition~2.1 of EXI-06-a), and hence neither cares about us nor is hostile to us (Theorem~3.1 of EXI-06-a), so that descriptions of it as hostile or pitiless are false (Corollary~3.2). We ask what this implies for meaning.

\section{Indifference does not remove meaning}

\begin{proposition}\label{p:meaning}
That the universe is indifferent to us does not imply that our lives lack meaning.
\end{proposition}

\begin{proof}[Argument]
On the account of EXI-01, a life is meaningful to the extent that its subject engages with projects of objective worth, whose value does not depend on the attitudes of those engaged in them (Definition~1 and Assumption~2 of EXI-01-a). On the account of EXI-17-a, which agrees, an object of engagement is worthwhile if its value does not depend on the agent's attitudes toward it (Definition~2.2 of EXI-17-a). If the worth of a project does not depend on the attitudes of those engaged in it, it does not depend on the attitudes of the universe either, and a fortiori not on the universe's lack of attitudes. The relief of suffering, the growth of understanding and the flourishing of persons have the worth they have whether or not anything in the cosmos registers it. Indifference, which is the absence of responsiveness, subtracts nothing.
\end{proof}

The argument does not deny that a caring universe would add something: it might add the good of being cared for. It denies that meaning in life requires that good.

\begin{proposition}\label{p:despair}
Cosmic indifference gives no ground for despair and no object for defiance.
\end{proposition}

\begin{proof}[Argument]
Despair over the universe presents it as having denied us something we needed; defiance presents it as an adversary. By Proposition~\ref{p:meaning}, its indifference denies us nothing that meaning requires; by Theorem~3.1 of EXI-06-a, it is not hostile and so is not an adversary. Both attitudes misrepresent their object. This agrees with the result of EXI-04-a that revolt is not a fitting response to the absurd, since the world it scorns is indifferent rather than opposed (Proposition~3.5 of EXI-04-a).
\end{proof}

Russell's ethic of defiance thus rests on the rhetoric that Corollary~3.2 of EXI-06-a rejects. What survives of it is the thought that our ideals are ours to maintain, and this needs no adversary: it follows from indifference alone, since nothing else will maintain them.

\section{Objections}

\begin{objection}
Meaning requires that our lives matter from the point of view of the universe. An indifferent universe has no such point of view, so nothing matters from it.
\end{objection}

\begin{reply}
If \enquote{the point of view of the universe} is the point of view of a subject that cares, then by Theorem~3.1 of EXI-06-a there is no such point of view, and the demand cannot be met by any universe without attitudes. If it is the impartial standpoint from which values are assessed without reference to any particular agent, then the worth of objects of engagement is precisely worth from that standpoint, and the demand is met \citep{Nagel1986}. On neither reading does indifference deprive our lives of meaning.
\end{reply}

\begin{objection}
Theistic accounts hold that meaning requires a purpose conferred by a being who cares for us. If the universe does not care, meaning is lost unless there is such a being.
\end{objection}

\begin{reply}
EXI-06-a is neutral on whether a being distinct from the universe cares for us, and so is this argument. The present point is conditional: if meaning consists in engagement with what has objective worth, as EXI-01 holds, then it does not require any being's care, the universe's or another's. Whether there is a further good of being cared for by such a being is a separate question, and supernaturalist accounts of meaning are assessed elsewhere \citep{Metz2013}.
\end{reply}

\section{Conclusion}

The universe does not care about us, and our lives may be meaningful all the same. The worth of what we engage with does not wait on cosmic recognition, and an indifferent world, being no enemy, calls neither for despair nor for defiance. It leaves the question of meaning where EXI-01 places it: in what we engage with, and how.

\begin{thebibliography}{9}\small
\bibitem[Metz(2013)]{Metz2013} Metz, T. (2013). \emph{Meaning in Life: An Analytic Study}. Oxford University Press.
\bibitem[Nagel(1986)]{Nagel1986} Nagel, T. (1986). \emph{The View from Nowhere}. Oxford University Press.
\bibitem[Russell(1903)]{Russell1903} Russell, B. (1903). A free man's worship. \emph{The Independent Review}, 1, 415--424.
\end{thebibliography}

\end{document}
